% ===========================================================================
\section{Sylow 定理}\label{sec:sylow}

\subsection{Sylow 理论：定义}

设 $G$ 为有限群，$|G| = p^k m$，$p$ 为素数且 $p\nmid m$（$p^k$ 是 $|G|$
中 $p$ 的最高幂次）。

\begin{definition}[$p$-子群]
$H\le G$ 称为 \keyword{$p$-子群}，如果它是子群且 $|H|$ 是 $p$ 的幂。
\end{definition}

\begin{definition}[Sylow $p$-子群]
$H\le G$ 称为 \keyword{Sylow $p$-子群}，如果它是 $p$-子群且
$|H| = p^k$——$p$ 的幂次取到最大。
\end{definition}

\begin{example}[$G = S_4$]
$|S_4| = 24 = 2^3\cdot 3$。
\begin{itemize}[leftmargin=1.3em,itemsep=1pt]
  \item $p=3$：Sylow $3$-子群大小为 $3$，例如 $\langle(1\,2\,3)\rangle$；
  \item $p=2$：Sylow $2$-子群大小应为 $8$；其存在性由下面的 Sylow 第一
        定理保证，具体构造见本节末尾的例子。
\end{itemize}
\end{example}

\subsection{Sylow 第一定理（存在性）}

\begin{theorem}[Sylow 第一定理]
对任何有限群 $G$ 与素数 $p$，Sylow $p$-子群存在。
\end{theorem}

\begin{proof}
（对 $|G|$ 归纳）设 $|G| = p^k m$，$p\nmid m$，$k\ge 1$。归纳假设：
命题对所有更小的群成立。考虑类方程
\[
  |G| = |C_G| + \sum_i [G : C_G(x_i)],
\]
分两种情形。

\keyword{情形 1：$p\nmid |C_G|$。}左边被 $p$ 整除而 $|C_G|$ 不被 $p$
整除，故右端至少有一项不被 $p$ 整除：存在 $x_i$ 使 $p\nmid [G:C_G(x_i)]$，
从而 $p^k \mid |C_G(x_i)|$。又 $C_G(x_i)$ 比 $G$ 小（因为
$[G:C_G(x_i)]>1$），由归纳假设其中有 $p^k$ 阶子群，即 $G$ 的 Sylow
$p$-子群。

\keyword{情形 2：$p\mid |C_G|$。}中心是阿贝尔群，由（阿贝尔群的）
Cauchy 定理，中心里有 $p$ 阶元素 $c$。考虑 $\langle c\rangle$：
\begin{itemize}[leftmargin=1.3em,itemsep=2pt]
  \item $\langle c\rangle\le C_G$，其中元素与 $G$ 的一切元素交换，
        故 $\langle c\rangle\normal G$；
  \item 商群 $G/\langle c\rangle$ 的阶为 $p^{k-1}m$，由归纳假设其中有
        $p^{k-1}$ 阶子群 $\bar H$；
  \item 取 $\bar H$ 在自然同态 $\pi: G\to G/\langle c\rangle$ 下的原像
        $H=\pi^{-1}(\bar H)$，则 $H\le G$ 且
        $|H| = |\bar H|\cdot|\langle c\rangle| = p^{k-1}\cdot p = p^k$，
        即为 Sylow $p$-子群。
\end{itemize}
\end{proof}

\begin{remark}[回到 $S_4$]
先看 $\{e,\ (1\,2),\ (3\,4),\ (1\,2)(3\,4)\}$——一个 4 阶 $2$-子群
（还不是 Sylow）；再把 $(1\,3)(2\,4)$ 加进来，膨胀成一个 8 阶子群——
这就是 $S_4$ 的 Sylow $2$-子群（同构于二面体群 $D_4$，阶为 $8$）。
\end{remark}

\subsection{Sylow 第二定理}

\begin{theorem}[Sylow 第二定理]
设 $P$ 是 $G$ 的任一 Sylow $p$-子群，$Q$ 是 $G$ 的任一 $p$-子群，
则存在 $g\in G$ 使得
\[
  Q \le gPg^{-1},
\]
即任何 $p$-子群都含于该 Sylow $p$-子群的某个共轭之中。
\end{theorem}

\begin{corollary}
\leavevmode
\begin{enumerate}[leftmargin=1.6em,itemsep=2pt]
  \item \keyword{任意两个 Sylow $p$-子群互相共轭。}取 $Q$ 也是 Sylow
        $p$-子群，则 $Q\le gPg^{-1}$，而两边大小相同，故 $Q = gPg^{-1}$。
  \item \keyword{任何 $p$-子群都含于某个 Sylow $p$-子群。}
  \item \keyword{Sylow $p$-子群唯一 $\iff$ 它是正规子群。}
        \begin{itemize}[leftmargin=1.3em,itemsep=1pt]
          \item $(\Leftarrow)$ 若 $P\normal G$：对任何 $g$，
                $gPg^{-1} = P$；而所有 Sylow $p$-子群都是 $P$ 的共轭，
                故只有自己一个。
          \item $(\Rightarrow)$ 若 $P$ 唯一：对任何 $g$，$gPg^{-1}$ 仍是
                Sylow $p$-子群（共轭不改变大小），由唯一性 $gPg^{-1}=P$，
                故 $P$ 正规。
        \end{itemize}
\end{enumerate}
\end{corollary}

\begin{proof}[Sylow 第二定理的证明]
设 $P$ 为 Sylow $p$-子群，$Q$ 为任意 $p$-子群，让 $Q$ \keyword{左乘作用}
在 $P$ 的全体左陪集组成的集合 $G/P$ 上：
\[
  q\cdot(gP) := (qg)P.
\]
注意这里的 $G/P$ \keyword{不是商群}，只是左陪集的集合（没有假设 $P$
正规）。不难验证这确实是一个群作用。
\begin{itemize}[leftmargin=1.3em,itemsep=2pt]
  \item 陪集个数 $[G:P] = m$ 不被 $p$ 整除（$P$ 已把 $|G|$ 中 $p$ 的
        幂次全用掉了）。而 $[G:P]$ 等于各轨道大小之和，所以
        \keyword{必存在一条轨道，其大小不被 $p$ 整除}；
  \item 另一方面，由轨道与稳定子定理，每条轨道大小整除 $|Q|$（$p$ 的幂）；
        不被 $p$ 整除的 $p$ 的幂只有 $1$——故存在\keyword{不动点} $gP$：
        对任意 $q\in Q$ 有 $qgP = gP$。
\end{itemize}
两边左乘 $g^{-1}$ 得 $g^{-1}qg\in P$，即
\[
  Q \subseteq gPg^{-1}.
\]
\end{proof}

\subsection{正规化子（Normalizer）}

\begin{definition}[正规化子（Normalizer）]
设 $P\le G$，$P$ 在 $G$ 中的\keyword{正规化子}为
\[
  N_G(P) = \{\,g\in G \mid gPg^{-1} = P\,\},
\]
即共轭作用下使 $P$ 不变的元素全体（也就是子群集合上共轭作用中 $P$ 的
稳定子）。
\end{definition}

\begin{itemize}[leftmargin=1.3em,itemsep=1pt]
  \item 易验证 $N_G(P)\le G$，且显然 $P\subseteq N_G(P)$；
  \item 顾名思义，$P\normal N_G(P)$——我们只挑选了使 $P$ 在共轭下不变
        的元素。
\end{itemize}

\begin{lemma}
设 $P$ 是 $G$ 的 Sylow $p$-子群，则 $P$ 是 $N_G(P)$ 的\keyword{唯一}
Sylow $p$-子群。
\end{lemma}

\begin{proof}
由 $P\le N_G(P)\le G$ 知 $|N_G(P)|$ 中 $p$ 的最高幂次仍为 $|P|$，故
$P$ 是 $N_G(P)$ 的 Sylow $p$-子群；又 $P\normal N_G(P)$，把上面推论
第 3 条（唯一 $\iff$ 正规）应用于群 $N_G(P)$，即得 $P$ 是 $N_G(P)$ 的
唯一 Sylow $p$-子群。
\end{proof}

\subsection{例子：对称群的 Sylow 子群的构造}

\begin{example}[$S_4$ 的 Sylow $2$-子群（阶 8）]
由如下置换生成：
\begin{itemize}[leftmargin=1.3em,itemsep=1pt]
  \item $\{1,2\}$ 与 $\{3,4\}$ 两块\keyword{整体交换}：$(1\,3)(2\,4)$；
  \item $\{1,2\}$ 内部交换：$(1\,2)$；
  \item $\{3,4\}$ 内部交换：$(3\,4)$。
\end{itemize}
\end{example}

\begin{example}[$S_{p^2}$ 的 Sylow $p$-子群]
类似想法可构造一般的 $S_n$ 的 Sylow $p$-子群。以 $S_{p^2}$ 为例：
\begin{itemize}[leftmargin=1.3em,itemsep=2pt]
  \item 把 $1,\dots,p^2$ 分成 $p$ 组：$\{1,\dots,p\}$、
        $\{p+1,\dots,2p\}$、…、$\{(p-1)p+1,\dots,p^2\}$；
  \item $p$ 个组之间可以\keyword{整体轮换}（1 号组→2 号组→…）；
  \item 每组内部各有一个 $p$-轮换。
\end{itemize}
这样得到一个 $p^{p+1}$ 阶子群，正是 $S_{p^2}$ 的 Sylow $p$-子群。
$n$ 更大时构造类似。
\end{example}

\begin{remark}
由 Sylow 第二定理，Sylow $p$-子群彼此共轭，所以其它方式构造出的
Sylow $p$-子群都与上述构造共轭。
\end{remark}

\subsection{Sylow 第三定理}

\begin{theorem}[Sylow 第三定理]
设 $|G| = p^n m$，$p\nmid m$，记
\[
  n_p = \#\{\,\text{$G$ 的 Sylow $p$-子群}\,\}.
\]
则
\[
  n_p \mid m,\qquad n_p \equiv 1 \pmod p.
\]
\end{theorem}

\begin{proof}[$n_p\mid m$ 的证明]
考虑 $G$ 在全体 Sylow $p$-子群组成的集合 $\Syl_p(G)$ 上的
\keyword{共轭作用}：
\begin{itemize}[leftmargin=1.3em,itemsep=2pt]
  \item 由 Sylow 第二定理，任意两个 Sylow $p$-子群共轭，故该作用
        \keyword{只有一个轨道}——这种作用称为\keyword{传递作用}
        （Transitive Action，亦称可迁）；
  \item 由轨道与稳定子定理，轨道大小 $n_p = [G : \Stab(P)]$，而这里的
        稳定子正是正规化子：
        \[
          \Stab(P) = \{\,g \mid gPg^{-1} = P\,\} = N_G(P);
        \]
  \item 故 $n_p = [G : N_G(P)]$。由于 $P\le N_G(P)$，
        $n_p = [G:N_G(P)]$ 整除 $[G:P] = m$。
\end{itemize}
\end{proof}

\begin{proof}[$n_p\equiv 1\pmod p$ 的证明]
取定一个 Sylow $p$-子群 $P$，让 $P$ \keyword{共轭作用}在 $\Syl_p(G)$
上：
\[
  g\cdot P_i := gP_i g^{-1}\qquad (g\in P).
\]
把 $n_p$ 按轨道分解，每个轨道取代表 $P_i$：
\[
  n_p = \sum |\Orb(P_i)| = \sum [\,P : \Stab(P_i)\,],
\]
其中 $\Stab(P_i) = \{\,g\in P \mid gP_ig^{-1} = P_i\,\} \le P$。
分两种情形：
\begin{itemize}[leftmargin=1.3em,itemsep=2pt]
  \item \keyword{$\Stab(P_i) \ne P$：}此时 $[P:\Stab(P_i)]$ 是大于 1
        的 $p$ 的幂，故 $p \mid |\Orb(P_i)|$；
  \item \keyword{$\Stab(P_i) = P$：}即 $P$ 中每个元素都使 $P_i$ 在共轭
        下不变，故 $P\le N_G(P_i)$。由上一小节的引理，$P_i$ 是
        $N_G(P_i)$ 的\keyword{唯一} Sylow $p$-子群；而
        $P\le N_G(P_i)$ 且 $|P| = |P_i|$，$P$ 也是 $N_G(P_i)$ 的
        Sylow $p$-子群，所以 $P = P_i$。这条轨道就是 $\{P\}$，
        大小为 $1$。
\end{itemize}
于是只有一条大小为 1 的轨道（$P$ 自己），其余轨道大小都被 $p$ 整除，故
\[
  n_p \equiv 1 \pmod p.
\]
\end{proof}
